\documentclass{article}
\usepackage[spanish]{babel}
\usepackage{arxiv}
\usepackage{enumitem}
\usepackage{url}
\usepackage[utf8]{inputenc} % allow utf-8 input
\usepackage[T1]{fontenc}    % use 8-bit T1 fonts
\usepackage{hyperref}       % hyperlinks
\usepackage{url}            % simple URL typesetting
\usepackage{booktabs}       % professional-quality tables
\usepackage{amsfonts}       % blackboard math symbols
\usepackage{nicefrac}       % compact symbols for 1/2, etc.
\usepackage{microtype}      % microtypography
\usepackage{lipsum}
\usepackage{graphicx}
\usepackage{float}
\graphicspath{ {./images/} }
\usepackage{float}
\usepackage{caption}
\usepackage{fancyhdr}
\pagestyle{fancy}
\fancyhf{} % limpia encabezado y pie
\fancyhead[R]{\textit{Constante-Ferrero-Gil}}
\fancyfoot[C]{\thepage}
\setlength{\textfloatsep}{10pt}
\setlength{\floatsep}{8pt}
\setlength{\intextsep}{8pt}
\setlength{\abovecaptionskip}{4pt}
\setlength{\belowcaptionskip}{4pt}
\title{TP 1.1 - SIMULACIÓN DE UNA RULETA}

\author{
 Constante Gonzalo \\
  Universidad Tecnologica Nacional\\
  Facultad Regional Rosario\\
Zeballos 1341, S2000, Argentina 
\\
  \texttt{constante.gonzalo2@gmail.com} \\
  %% examples of more authors
   \And
 Gil Francisco Marcelo\\
  Universidad Tecnologica Nacional\\
  Facultad Regional Rosario\\
Zeballos 1341, S2000, Argentina 
 \\
  \texttt{fran17mg@gmail.com} \\
  \And
 Ferrero Santiago\\
  Universidad Tecnologica Nacional\\
  Facultad Regional Rosario\\
Zeballos 1341, S2000, Argentina 
\\
  \texttt{santifnob@gmail.com} \\
}

\begin{document}
\maketitle
\begin{abstract}
El presente trabajo práctico tiene como finalidad desarrollar y analizar un modelo computacional de simulación de una ruleta europea, utilizando herramientas de programación en Python y conceptos fundamentales de probabilidad y estadística. La ruleta considerada posee 37 posibles resultados, numerados del 0 al 36, todos ellos asumidos como equiprobables bajo un comportamiento ideal del sistema. A partir de esta hipótesis, se implementó un programa capaz de generar múltiples tiradas aleatorias y organizar los resultados en distintas corridas independientes, con el objetivo de estudiar el comportamiento experimental de variables estadísticas relevantes y compararlas con sus valores teóricos esperados.

Para llevar adelante la simulación se emplearon estructuras básicas de programación, tales como listas, ciclos iterativos y funciones, junto con bibliotecas específicas para el cálculo numérico y la visualización de datos. El programa recibe parámetros por consola, entre ellos el número elegido, la cantidad de tiradas y la cantidad de corridas, lo cual permite modificar las condiciones del experimento sin alterar la estructura principal del código.

El análisis realizado se centra principalmente en la frecuencia relativa de aparición de un número seleccionado, el valor promedio de las tiradas, la varianza, el desvío estándar y la distribución de los valores generados. A nivel teórico, se espera que la frecuencia relativa de cualquier número tienda a $\frac{1}{37}$ conforme aumenta la cantidad de tiradas, ya que cada resultado posee la misma probabilidad de ocurrencia. Del mismo modo, el valor promedio de los números obtenidos debería aproximarse a 18, mientras que la varianza debería tender a 114 y el desvío estándar a un valor cercano a 10,67.

A partir de las gráficas generadas, se espera observar que en corridas con pocas tiradas los resultados presentan oscilaciones importantes, propias del azar y del tamaño reducido de la muestra. Sin embargo, al aumentar la cantidad de tiradas y al promediar varias corridas, dichas fluctuaciones deberían disminuir progresivamente, mostrando una tendencia hacia los valores teóricos. Este comportamiento permite vincular la simulación con conceptos como la ley de los grandes números, la estabilidad de las frecuencias relativas y el análisis de medidas muestrales. El trabajo no solo busca reproducir computacionalmente el funcionamiento de una ruleta, sino también utilizarla como caso de estudio para comprender cómo los fenómenos aleatorios pueden ser modelados, observados y contrastados mediante herramientas estadísticas.
\end{abstract}

\keywords{Simulación \and Ruleta \and Python \and Probabilidad \and Estadística \and Frecuencia relativa \and Varianza}


% keywords can be removed
%\keywords{First keyword \and Second keyword \and More}

\newpage

\section{Introducción}

La simulación computacional constituye una herramienta fundamental para el estudio de fenómenos aleatorios y sistemas complejos, en los cuales el análisis analítico directo puede resultar limitado o impracticable. En este contexto, la estadística y la teoría de probabilidades permiten describir y modelar el comportamiento esperado de dichos sistemas a través de variables aleatorias y distribuciones teóricas.

En el presente trabajo práctico se estudia el comportamiento de una ruleta europea mediante simulación computacional. Este sistema, aunque asociado a un juego de azar, resulta especialmente útil desde el punto de vista académico, ya que permite analizar de manera empírica conceptos fundamentales como la probabilidad, la frecuencia relativa, la media muestral y la variabilidad de los resultados en experimentos aleatorios repetidos.

La ruleta europea está compuesta por 37 resultados posibles, numerados del 0 al 36, y en un modelo ideal cada uno de estos valores posee la misma probabilidad de ocurrencia. Esto permite representar el experimento mediante una variable aleatoria discreta con distribución uniforme:

\begin{equation}
    P(X = x) = \frac{1}{37}, \quad x \in \{0,1,2,\dots,36\}
\end{equation}

A partir de este modelo teórico, el objetivo del trabajo consiste en analizar cómo se comportan empíricamente distintas medidas estadísticas cuando el experimento es repetido un gran número de veces. En particular, se busca estudiar la evolución de la media, la frecuencia relativa de eventos específicos y la dispersión de los resultados, comparando el comportamiento observado con los valores teóricos esperados.

Asimismo, se pretende evaluar cómo el aumento en la cantidad de observaciones influye en la estabilidad de los estimadores muestrales, permitiendo observar fenómenos fundamentales de la teoría de la probabilidad como la Ley de los Grandes Números y la tendencia hacia la convergencia de los estadísticos.

Finalmente, el análisis de los resultados obtenidos mediante simulación permite establecer una conexión entre el modelo teórico y su comportamiento experimental, destacando la utilidad de la simulación computacional como herramienta para el estudio y comprensión de procesos aleatorios.

\newpage

\section{Materiales y Métodos}

\subsection{Conceptos teóricos empleados}

Para el estudio de dicho juego de azar, vamos a utilizar diferentes conceptos importantes del campo de la Probabilidad y Estadística, que nos van a ayudar a predecir el comportamiento de la ruleta y estudiar los resultados provistos de la simulación.
De antemano necesitamos entender estos conceptos, para luego comprender el análisis de los resultados. Dentro de estos conceptos encontramos: 
\begin{itemize}
    \item \textbf{Población}: Conjunto total, completo y bien definido de individuos, objetos, eventos o medidas que comparten características comunes y sobre los cuales se desea realizar un estudio o inferencia.
    
    \item \textbf{Muestra}: Subconjunto representativo y finito que se extrae de un conjunto mayor, llamado población, la cual puede ser finita o infinita. Permite inferir resultados de la población sin necesidad de analizar a todos los elementos.
    
    \item \textbf{Estadística descriptiva}: Técnicas para resumir y visualizar datos con gráficos y medidas.
    
    \item  \textbf{Estadística Inferencial}: Métodos para deducir propiedades de una población a partir de una 
    muestra, permitiendo realizar predicciones (estimaciones, contrastes de hipótesis). 
    \item \textbf{Probabilidad}: Grado de certeza de que un evento ocurra, expresado en un número entre 0 y 1. 
    
    \item \textbf{Distribuciones de Probabilidad}: Modelos teóricos que describen cómo se espera que varíen los datos 
    (ej. Distribución Normal o Binomial).  Se modelan, definen y caracterizan en base a parámetros predefinidos. 
    
    \item \textbf{Variable Aleatoria}: función matemática que asigna un valor numérico específico a cada resultado posible de un experimento aleatorio. Se clasifican principalmente en discretas (valores contables) y continuas (valores en un intervalo). 
    

\end{itemize}
En cuanto a las \textbf{medidas} más importantes para este estudio,  que nos permiten cuantificar y estudiar los datos obtenidos, vemos:
\begin{itemize}
    \item \textbf{Frecuencia absoluta}: Número total de veces que un valor, dato o evento aparece en un experimento o estudio. 
    \item \textbf{Frecuencia relativa}: Cociente (división) entre la frecuencia absoluta  y el tamaño total de la muestra. 
\end{itemize}
\begin{itemize}

    \item \textbf{Media Aritmética Poblacional:} Promedio teórico de todos los valores de una población completa.
    \[ \mu = \frac{1}{N} \sum_{i=1}^{N} X_i \]
    
    \item \textbf{Media Aritmética Muestral:} Promedio calculado a partir de un subconjunto de datos extraído de la población.
    \[ \bar{x} = \frac{1}{n} \sum_{i=1}^{n} x_i \]
    
    \item \textbf{Varianza Poblacional:} Promedio del cuadrado de las desviaciones de todos los elementos de la población respecto a la media poblacional.
    \[ \sigma^2 = \frac{1}{N} \sum_{i=1}^{N} (X_i - \mu)^2 \]
    
    \item \textbf{Varianza Muestral (Insesgada):} Estimador de la varianza poblacional calculado a partir de una muestra, corregido con el factor de Bessel ($n-1$).
    \[ s^2 = \frac{1}{n-1} \sum_{i=1}^{n} (x_i - \bar{x})^2 \]
    
    \item \textbf{Desviación Estándar:} Raíz cuadrada de la varianza. Expresa la dispersión en las mismas unidades originales de la variable estudiada.
    \[ \sigma = \sqrt{\sigma^2} \quad \text{ó} \quad s = \sqrt{s^2} \]

\end{itemize}

Luego encontramos los diferentes \textbf{parámetros}, que entendemos por ellos como valores aleatorios e invariables que definen formalmente el comportamiento global de una variable aleatoria $X$. 

\begin{itemize}
    \item \textbf{Esperanza Matemática (Caso Continuo):} El centro de gravedad o valor esperado a largo plazo de una variable aleatoria continua ponderada por su densidad de probabilidad.
    \[ E[X] = \int_{-\infty}^{\infty} x \cdot f(x) \, dx \]
    
    \item \textbf{Esperanza Matemática (Caso Discreto):} Suma ponderada de todos los valores posibles de una variable aleatoria discreta multiplicados por su respectiva probabilidad puntual.
    \[ E[X] = \sum_{i} x_i \cdot P(X = x_i) \]
    
    \item \textbf{Varianza de una Variable Aleatoria:} Esperanza del cuadrado de la distancia de la variable con respecto a su media teórica. Mide la variabilidad aleatoria.
    \[ Var(X) = E[(X - E[X])^2] = E[X^2] - (E[X])^2 \]
\end{itemize}

Y por ultimo, las \textbf{distribuciones} que vamos a tener que aplicar van a ser las siguientes:
\begin{itemize}
    \item \textbf{Distribución Normal:} $X \sim N(\mu, \sigma^2)$. Modeliza variables naturales continuas. Parámetros: media o centro de la campana ($\mu \in \mathbb{R}$) y la varianza o ancho de la misma ($\sigma^2 > 0$).
    \item \textbf{Distribución Uniforme Discreta:} $X \sim U(a, b)$. Modeliza situaciones con un número finito de resultados posibles y aislados (como el lanzamiento de un dado), donde cada valor posee exactamente la misma probabilidad de ocurrir. Parámetros: valor mínimo entero ($a \in \mathbb{Z}$) y valor máximo entero ($b \in \mathbb{Z}$ con $b \ge a$), donde el número total de resultados es $n = b - a + 1$.
\end{itemize}

La variable aleatoria \textbf{$X$}, establecida como lanzamientos de la ruleta, posee una \textbf{distribución uniforme discreta} como se menciono en la introducción, la cual representamos por: 

\begin{center}    \textbf{$X \sim U(0, 36)$}
\end{center}

La probabilidad de que salga cada valor es de aproximadamente 0.27, dato que nos va a servir como la \textbf{frecuencia relativa esperada} posteriormente, cuando estudiemos el caso de dicha frecuencia para cierto numero dado con respecto a multiples corridas de n tiradas.

Además, a partir de esta distribución de $X$, podemos calcular la esperanza matemática, la varianza, y su desviación estándar, medidas las cuales vamos a usar más adelante como \textbf{valores esperados}. El calculo se lleva a cabo de la siguiente manera:

\begin{equation}
    \mu = E(X) = \frac{a + b}{2} = \frac{0 + 36}{2} = \frac{36}{2} = 18
    \label{eq:media_uniforme}
\end{equation}

\begin{equation}
    \sigma^{2} = \operatorname{Var}(X) = \frac{(b - a + 1)^{2} - 1}{12} = \frac{(36 - 0 + 1)^{2} - 1}{12} = \frac{37^{2} - 1}{12} = \frac{1369 - 1}{12} = \frac{1368}{12} = 114
    \label{eq:varianza_uniforme}
\end{equation}

\begin{equation}
    \sigma = \sqrt{\operatorname{Var}(X)} = \sqrt{114} \approx 10.677
    \label{eq:desvio_uniforme}
\end{equation}

Sobre algo en lo que recaemos mucho a lo largo de este informe, es la \textbf{ley de los grandes números}, la cual establece que, al aumentar el tamaño de la muestra, los estadísticos muestrales tienden a \textbf{aproximarse a los parámetros teóricos de la población}. En este trabajo, este principio permite justificar que la frecuencia relativa y las medidas experimentales converjan hacia los valores esperados de la ruleta al incrementar el número de tiradas.

Y respecto al concepto anterior, el mismo esta relacionado directamente con el \textbf{teorema central del limite (TCL)}, el cual establece que, al tomar \textbf{muestras independientes y suficientemente grandes} (N >= 30) de cualquier población (con media y varianza finitas), la distribución de sus medias muestrales tenderá a una \textbf{distribución normal}, incluso si la población original no sigue una distribución normal. Se utilizará para demostrar que las medias muestrales de las corridas en la ruleta van a \textbf{tender a una distribución normal}, siendo que la población de estudio (las tiradas) se distribuyen inicialmente con una distribución uniforme discreta.

\subsection{Implementación computacional}

Para la realización del programa de simulación del plato de una ruleta, hemos utilizado el lenguaje de programación \textbf{Python 3.x}, utilizando paquetes que facilitan desde la obtención de datos al azar, hasta la creación de diferentes gráficas. Estos son:

\begin{itemize}
    \item \textbf{Random}: Módulo de la biblioteca estándar utilizado para generar números pseudoaleatorios y realizar selecciones aleatorias. Ideal para este tipo de simulaciones, ya que permite crear números enteros, barajar secuencias y seleccionar elementos al azar de forma rápida y sencilla. Utiliza el algoritmo \textbf{Mersenne Twister} como núcleo para la generación de números pseudoaleatorios, el cual posee un período lo suficientemente largo ($2^{19937}$) para garantizar la independencia estadística requerida en este experimento.
    \item \textbf{Matplotlib}: biblioteca integral de Python para la creación de visualizaciones estáticas, animadas e interactivas. En nuestro caso, nos facilitará la creación de los gráficos que mostrarán las fluctuaciones de las medidas, como la frecuencia relativa, a través de las tiradas de las corridas.
    \item \textbf{Numpy}: Biblioteca fundamental para la computación científica. La utilizaremos para realizar cálculos como la media, varianza y desvío estándar sobre los datos aleatorios.
    \item \textbf{Sys}: Este módulo proporciona funciones y variables específicas del sistema para interactuar con el intérprete de Python y el entorno operativo. Nos permitirá introducir los parámetros necesarios a la hora de correr la simulación, como el número elegido para medir su frecuencia relativa, las corridas, y la cantidad de tiradas por cada corrida.
\end{itemize}

Para poder correr el programa, hay que pararse sobre el directorio que lo contiene, y ejecutar el siguiente comando en consola:
\begin{verbatim}
$ python main.py -n <num_valores> <tiradas> <corridas>
\end{verbatim}
Donde cada parámetro se reemplaza por:
\begin{enumerate}
    \item \textbf{<num\_valores>} -> número, entre el 0 y 36, utilizado para el estudio del comportamiento frecuencia relativa.
    \item \textbf{<corridas>} -> cantidad de corridas a realizar
    \item \textbf{<tiradas>} -> cantidad de tiradas a realizar por cada corrida
\end{enumerate}

Una vez ejecutado, va a llevar a cabo una secuencia de pasos para poder realizar la simulación de forma correcta. Estos son: 
\begin{enumerate}
    \item Tomar los valores de los parámetros ingresados por consola con el array de argumentos que provee sys (\textbf{sys.argv[numero\_argumento]}).
    \item Se llama la función \textbf{generar\_valores\_aleatorios(tiradas, corridas)}, la cual utiliza \textbf{random.randint(0,36)} y dos iteraciones for para formar un conjunto de datos (array), el cual va a ser bidimensional, donde la fila representa la corrida, y cada columna una tirada. Este array es retornado para poder ser usado en los diferentes cálculos.
    \item Para el análisis de las medidas principales (\textit{frecuencia relativa, valor promedio, varianza y desvío estándar}), el programa sigue una lógica de visualización en tres partes mediante la librería \textbf{matplotlib}. En cada una de estas funciones se crea una figura con tres paneles (\textbf{subplots(1, 3)}) para contrastar la evolución de los datos:\begin{itemize}
    \item \textbf{Gráfico 1 (Individual):} Muestra la evolución de la métrica únicamente para la primera corrida, permitiendo observar la volatilidad inicial del azar.
    \item \textbf{Gráfico 2 (Múltiples corridas):} Superpone las curvas de todas las corridas ejecutadas, utilizando transparencias (\textbf{alpha=0.7}) para visualizar la dispersión entre diferentes experimentos independientes.
    \item \textbf{Gráfico 3 (Promedio de corridas):} Representa el promedio de los valores obtenidos en todas las corridas para cada tirada. Este gráfico se diseña con el fin de analizar la estabilidad del estimador hacia el valor teórico esperado.\end{itemize}En todos los casos, se traza una línea horizontal (\textbf{axhline}) que representa el valor teórico esperado (por ejemplo, 18 para la media o 1/37 para la frecuencia), facilitando la comparación visual entre lo empírico y lo teórico.
    \item En el cálculo de la dispersión, se utilizan funciones de la librería \textbf{numpy}, como \textbf{np.var()} para la varianza y \textbf{np.std()} para el desvío estándar. Estos cálculos se realizan de forma acumulativa (tirada a tirada) para observar cómo el estimador se estabiliza a medida que aumenta el tamaño de la muestra ($n$).
    \item Como etapa final del procesamiento de datos, se llama a la función \textbf{graficar\_histogramas(resultados, tiradas)}. Esta función genera dos tipos de visualizaciones fundamentales utilizando \textbf{plt.hist()}:\begin{itemize}
    \item \textbf{Distribución de valores individuales:} Se aplanan todos los datos generados para comprobar que la frecuencia de aparición de cada número del 0 al 36 tiende a ser uniforme.
    \item \textbf{Distribución de medias muestrales:} Se obtiene tomando el valor final $\bar{X}$ de cada una de las $M$ corridas. Según el \textbf{Teorema Central del Límite}, a pesar de que la fuente original es uniforme, la distribución de estas medias debe aproximarse a una normal $N(\mu, \sigma^2/n)$.\end{itemize}
    \item Una vez generados todos los cálculos y procesadas las imágenes, el programa utiliza el método \textbf{savefig()} para exportar cada conjunto de gráficos como archivos de imagen independientes (.png), lo que facilita su posterior inclusión en este informe, y finalmente muestra las ventanas en pantalla mediante \textbf{plt.show()}.
    
\end{enumerate}

\newpage

\section{Resultados}

En esta sección se presentan y analizan los resultados obtenidos a partir de la simulación computacional de una ruleta europea. El programa desarrollado permite generar tiradas aleatorias dentro del intervalo discreto $[0,36]$, organizar los resultados en distintas corridas y calcular estadísticos relevantes como la frecuencia relativa, el valor promedio, la varianza y el desvío estándar. Además, se incorporan histogramas para estudiar tanto la distribución de los valores individuales como el comportamiento de las medias muestrales.

\subsection{Frecuencia relativa}

La frecuencia relativa esperada para la aparición de un número específico en una ruleta europea es:

\begin{equation}
    f_e = \frac{1}{37} \approx 0{,}027.
\end{equation}

Este valor surge de considerar que la ruleta posee 37 resultados posibles, numerados del 0 al 36, y que todos ellos tienen la misma probabilidad de ocurrencia. En consecuencia, para cualquier número seleccionado, la probabilidad teórica de aparición en una tirada aislada es $1/37$.

\begin{figure}[H]
    \centering
    \includegraphics[width=0.85\textwidth]{figura_01_frecuencia_relativa_corrida_1.jpg}
    \caption{Frecuencia relativa del número elegido en la primera corrida.}
    \label{fig:frecuencia-relativa-corrida-1}
\end{figure}

En la Figura \ref{fig:frecuencia-relativa-corrida-1} se observa una curva con oscilaciones marcadas durante las primeras tiradas. En el tramo inicial no se obtuvo el valor ingresado como número elegido, por lo que la frecuencia relativa permanece en cero. Posteriormente, se registra un incremento pronunciado que alcanza aproximadamente el valor $0{,}040$. Luego, al superar las 400 tiradas, la curva comienza a aproximarse al valor de frecuencia esperada, aunque todavía con oscilaciones propias del comportamiento aleatorio del experimento.

\begin{figure}[H]
    \centering
    \includegraphics[width=0.85\textwidth]{figura_02_frecuencia_relativa_6_corridas.jpg}
    \caption{Frecuencia relativa del número elegido en seis corridas independientes.}
    \label{fig:frecuencia-relativa-6-corridas}
\end{figure}

La Figura \ref{fig:frecuencia-relativa-6-corridas} compara la frecuencia relativa obtenida en seis corridas independientes. La leyenda del gráfico permite identificar el color correspondiente a cada corrida. En términos generales, todas las curvas presentan fuertes fluctuaciones iniciales y luego adoptan un comportamiento similar, con una tendencia progresiva hacia la frecuencia esperada. Se destaca que la corrida 5 registra un pico cercano a $0{,}5$ durante las primeras tiradas, debido a que el número seleccionado apareció en una etapa temprana del experimento. Sin embargo, a medida que aumenta el número de tiradas, dicha curva pierde ese comportamiento extremo y comienza a asemejarse al resto.


\begin{figure}[H]
    \centering
    \includegraphics[width=0.85\textwidth]{figura_03_frecuencia_relativa_promediada_6_corridas.jpg}
    \caption{Frecuencia relativa promediada entre seis corridas.}
    \label{fig:frecuencia-relativa-promediada}
\end{figure}

La Figura \ref{fig:frecuencia-relativa-promediada} muestra la frecuencia relativa promediada entre las seis corridas. Esta curva resume el comportamiento conjunto de las simulaciones y permite reducir el ruido generado por la variabilidad individual de cada experimento. Se observa un pico inicial cercano a $0{,}1$, influenciado principalmente por el valor elevado registrado en la corrida 5. No obstante, aproximadamente a partir de la tirada 70, la curva comienza a estabilizarse alrededor del valor teórico esperado.


\subsection{Valor promedio}

El valor promedio muestral se calcula como la suma acumulada de los resultados obtenidos dividida por la cantidad de tiradas realizadas hasta cada instante. Para una ruleta europea ideal, el valor esperado de una tirada es:

\begin{equation}
    E(X) = \frac{0+1+2+\dots+36}{37} = \frac{666}{37}=18.
\end{equation}

Por lo tanto, se espera que el promedio muestral tienda progresivamente al valor 18 conforme aumenta la cantidad de tiradas.


\begin{figure}[H]
    \centering
    \includegraphics[width=0.85\textwidth]{figura_04_valor_promedio_corrida_1.jpg}
    \caption{Valor promedio acumulado de las tiradas en la primera corrida.}
    \label{fig:valor-promedio-corrida-1}
\end{figure}

En la Figura \ref{fig:valor-promedio-corrida-1}, la curva inicia en 36, lo cual indica que el primer resultado de la corrida fue justamente ese valor. En las primeras iteraciones se observan fluctuaciones importantes, dado que cada nuevo resultado tiene un peso considerable sobre el promedio acumulado. Sin embargo, al superar aproximadamente las 600 tiradas, la curva comienza a converger de manera asintótica hacia el promedio teórico esperado de 18.


\begin{figure}[H]
    \centering
    \includegraphics[width=0.85\textwidth]{figura_05_valor_promedio_6_corridas.jpg}
    \caption{Valor promedio acumulado de cada una de las seis corridas.}
    \label{fig:valor-promedio-6-corridas}
\end{figure}

La Figura \ref{fig:valor-promedio-6-corridas} presenta el valor promedio acumulado de cada una de las seis corridas. Cada curva posee un color distinto para facilitar su comparación. Se evidencia una dispersión inicial considerable, ya que en la primera tirada los valores pueden ubicarse muy por encima o por debajo del promedio esperado. A medida que el número de iteraciones aumenta, esta dispersión disminuye y las curvas comienzan a estabilizarse en valores cercanos a 18.


\begin{figure}[H]
    \centering
    \includegraphics[width=0.85\textwidth]{figura_06_valor_promedio_promediado_6_corridas.jpg}
    \caption{Valor promedio promediado entre seis corridas.}
    \label{fig:valor-promedio-promediado}
\end{figure}

En la Figura \ref{fig:valor-promedio-promediado} se representa el promedio de los valores promedio acumulados de las seis corridas. La curva inicia con un valor superior a $21{,}5$ y luego comienza a aproximarse al valor teórico. Alrededor de la tirada 50 se observa una primera estabilización, mientras que luego de aproximadamente 800 tiradas la curva converge con mayor claridad hacia el promedio esperado de 18.


\subsection{Varianza}

La varianza permite medir la dispersión de los resultados respecto del promedio. Para una variable aleatoria discreta uniforme con 37 valores consecutivos, la varianza teórica poblacional se calcula como:

\begin{equation}
    Var(X)=\frac{n^2-1}{12}=\frac{37^2-1}{12}=114.
\end{equation}

Este valor funciona como referencia para evaluar el comportamiento experimental de la varianza obtenida en la simulación.

\begin{figure}[H]
    \centering
    \includegraphics[width=0.85\textwidth]{figura_07_varianza_corrida_1.jpg}
    \caption{Varianza muestral acumulada de las tiradas en la primera corrida.}
    \label{fig:varianza-corrida-1}
\end{figure}

La Figura \ref{fig:varianza-corrida-1} muestra la varianza empírica correspondiente a la primera corrida. Durante las primeras tiradas se registra un salto abrupto que supera el valor 120. Luego, en tiradas posteriores, la curva oscila aproximadamente entre 100 y 120. En este tramo se mantiene levemente por debajo de la varianza esperada, aunque la amplitud de las fluctuaciones disminuye conforme aumenta el número de iteraciones.


\begin{figure}[H]
    \centering
    \includegraphics[width=0.85\textwidth]{figura_08_varianza_6_corridas.jpg}
    \caption{Varianza muestral acumulada de cada una de las seis corridas.}
    \label{fig:varianza-6-corridas}
\end{figure}

En la Figura \ref{fig:varianza-6-corridas} se observa la varianza acumulada de cada una de las seis corridas. Al igual que en los gráficos anteriores, la dispersión inicial es elevada. Algunas curvas registran valores muy alejados de la varianza teórica, como la corrida 3, que supera el valor 250, o la corrida 6, que inicia con varianza nula. No obstante, al superar aproximadamente las 200 iteraciones, las curvas comienzan a estabilizarse en valores próximos a 114. A partir de la iteración 400, la mayoría se mantiene levemente por debajo del valor esperado, con excepción de la corrida 5.


\begin{figure}[H]
    \centering
    \includegraphics[width=0.85\textwidth]{figura_09_varianza_promediada_6_corridas.jpg}
    \caption{Varianza muestral promediada entre seis corridas.}
    \label{fig:varianza-promediada}
\end{figure}

La Figura \ref{fig:varianza-promediada} expone el comportamiento de la varianza al promediar los valores muestrales de las seis corridas simultáneas. Durante las tiradas iniciales, el valor asciende rápidamente y supera los 100. Luego, la curva presenta una tendencia de aproximación asintótica hacia la varianza teórica de 114, con oscilaciones cada vez menos pronunciadas.


\subsection{Desvío estándar}

El desvío estándar se obtiene como la raíz cuadrada de la varianza. Para la ruleta simulada, su valor teórico esperado es:

\begin{equation}
    \sigma = \sqrt{114} \approx 10{,}67.
\end{equation}

Este estadístico permite interpretar la dispersión de los resultados en las mismas unidades que los valores originales de la ruleta.

\begin{figure}[H]
    \centering
    \includegraphics[width=0.85\textwidth]{figura_10_desvio_estandar_corrida_1.jpg}
    \caption{Desvío estándar acumulado de las tiradas en la primera corrida.}
    \label{fig:desvio-corrida-1}
\end{figure}

La Figura \ref{fig:desvio-corrida-1} muestra la evolución del desvío estándar en la primera corrida. En las tiradas iniciales se registra un pico superior a 11. A partir de aproximadamente la tirada 50, los valores oscilan entre 10 y el valor esperado. Luego de las 400 tiradas, la amplitud de las oscilaciones disminuye y el valor se mantiene próximo al desvío estándar teórico, aunque levemente por debajo.


\begin{figure}[H]
    \centering
    \includegraphics[width=0.85\textwidth]{figura_11_desvio_estandar_6_corridas.jpg}
    \caption{Desvío estándar acumulado de cada una de las seis corridas.}
    \label{fig:desvio-6-corridas}
\end{figure}

La Figura \ref{fig:desvio-6-corridas} presenta el desvío estándar de cada una de las seis corridas. En las primeras tiradas se observa una dispersión importante entre las curvas, producto del bajo tamaño muestral. Sin embargo, al aumentar la cantidad de iteraciones, las curvas comienzan a agruparse y a converger hacia el valor esperado de $10{,}67$, aproximadamente desde la tirada 50.


\begin{figure}[H]
    \centering
    \includegraphics[width=0.85\textwidth]{figura_12_desvio_estandar_promediado_6_corridas.jpg}
    \caption{Desvío estándar promediado entre seis corridas.}
    \label{fig:desvio-promediado}
\end{figure}

En la Figura \ref{fig:desvio-promediado} se representa el desvío estándar promediado de las seis corridas. La curva inicia en cero, debido a que al comienzo no existe dispersión acumulada suficiente, y luego presenta un incremento pronunciado hasta alcanzar un valor cercano a 10. A partir de la iteración 200, las oscilaciones se atenúan y la curva tiende progresivamente hacia el valor esperado de $10{,}67$.


\subsection{Histogramas}

Además de los gráficos acumulativos, se utilizaron histogramas para analizar la distribución general de los resultados obtenidos en la simulación.

\begin{figure}[H]
    \centering
    \includegraphics[width=0.85\textwidth]{figura_13_histograma_distribucion_uniforme.jpg}
    \caption{Histograma de frecuencia de aparición de los números de la ruleta.}
    \label{fig:histograma-uniforme}
\end{figure}

La Figura \ref{fig:histograma-uniforme} muestra la frecuencia de aparición de cada número individual, desde el 0 hasta el 36. Las barras presentan alturas similares y se ubican alrededor de la probabilidad teórica $1/37 \approx 0{,}027$. Este comportamiento confirma empíricamente que los números tienen aproximadamente la misma probabilidad de salir, lo cual se corresponde con una variable aleatoria de distribución uniforme discreta.


\begin{figure}[H]
    \centering
    \includegraphics[width=0.85\textwidth]{figura_14_histograma_medias_muestrales.jpg}
    \caption{Histograma de densidad de las medias muestrales obtenidas en la simulación.}
    \label{fig:histograma-medias}
\end{figure}

La Figura \ref{fig:histograma-medias} presenta un histograma de densidad correspondiente a las medias muestrales obtenidas luego de procesar 20000 tiradas. Se observa una morfología aproximadamente simétrica y acampanada, con un pico máximo cercano al valor medio teórico de 18. Hacia los extremos alejados del centro, la densidad disminuye de forma progresiva hasta alcanzar valores casi nulos. Este comportamiento se vincula con el Teorema Central del Límite, dado que la distribución de medias muestrales tiende a aproximarse a una distribución normal cuando el tamaño de la muestra es suficientemente grande.

\newpage

\section{Discusión}

El análisis de las gráficas permite comparar el comportamiento de los resultados obtenidos mediante el programa desarrollado en Python con los valores teóricos esperados desde la probabilidad y la estadística. En primer lugar, las curvas correspondientes a la frecuencia relativa, el valor promedio, la varianza y el desvío estándar exhiben un comportamiento inicial disperso y con oscilaciones marcadas. Esta situación es esperable en muestras pequeñas, ya que cada nuevo resultado tiene un impacto considerable sobre los estadísticos acumulados.

A medida que la cantidad de tiradas aumenta, se observa que todas las curvas atenúan la amplitud de sus oscilaciones y tienden de manera asintótica hacia sus respectivos valores teóricos. Este comportamiento constituye una evidencia empírica de la Ley de los Grandes Números, según la cual, al aumentar el tamaño de la muestra, las medidas experimentales tienden a aproximarse a los parámetros poblacionales.

En cuanto a las curvas obtenidas mediante el promedio de seis corridas, se observa una reducción significativa del ruido visual respecto de las curvas individuales. Al promediar los valores de varias corridas, se obtiene una curva más suave y estable. Esto permite visualizar con mayor claridad la tendencia general del sistema y demuestra que el estadístico promediado alcanza una estabilidad más rápida que una corrida individual aislada.

Por otro lado, para el análisis de los histogramas se aumentó la cantidad de tiradas con el objetivo de obtener una representación más clara del comportamiento probabilístico. El histograma de frecuencias individuales confirmó empíricamente que la ruleta simulada se comporta como una variable aleatoria de distribución uniforme discreta. La similitud en la altura de las barras permite corroborar que el programa asigna, de manera efectiva, probabilidades aproximadamente iguales a todos los elementos del espacio muestral.

El segundo histograma, correspondiente a la distribución de las medias muestrales, presenta una forma compatible con una distribución normal o campana de Gauss. Esta observación permite vincular el resultado con el Teorema Central del Límite, ya que las medias muestrales se concentran alrededor de la media poblacional teórica de 18. Además, al considerar el desvío estándar de las medias, se puede aplicar la regla empírica para analizar la proporción de valores esperada dentro de distintos intervalos alrededor de la media.

Si se considera un desvío estándar de las medias igual a $5{,}339$, los intervalos asociados a la regla empírica quedan expresados de la siguiente manera:

\begin{itemize}
    \item Aproximadamente el 68\% de los promedios obtenidos se encuentra dentro de un desvío estándar respecto de la media:
    \begin{equation}
        18 - 5{,}339 = 12{,}661 \quad \text{y} \quad 18 + 5{,}339 = 23{,}339.
    \end{equation}
    Por lo tanto, el intervalo aproximado es $[12{,}7; 23{,}3]$.

    \item Aproximadamente el 95\% de los datos se ubica dentro de dos desvíos estándar:
    \begin{equation}
        18 - 2(5{,}339) = 7{,}322 \quad \text{y} \quad 18 + 2(5{,}339) = 28{,}678.
    \end{equation}
    Por lo tanto, el intervalo aproximado es $[7{,}3; 28{,}7]$.

    \item Aproximadamente el 99{,}7\% de los resultados se encuentra dentro de tres desvíos estándar:
    \begin{equation}
        18 - 3(5{,}339) = 1{,}983 \quad \text{y} \quad 18 + 3(5{,}339) = 34{,}017.
    \end{equation}
    Por lo tanto, el intervalo aproximado es $[2{,}0; 34{,}0]$.
\end{itemize}

\newpage

\section{Conclusión}

La implementación de la aplicación en Python permitió gestionar de manera eficiente las variables de entrada, tales como el número seleccionado, la cantidad de tiradas y el número de corridas. Asimismo, posibilitó una representación gráfica detallada de los resultados obtenidos. El empleo de bibliotecas especializadas resultó fundamental para el desarrollo del trabajo: \texttt{Matplotlib} facilitó la visualización de datos, el módulo \texttt{random} permitió la generación de valores aleatorios en el intervalo discreto $[0,36]$ y \texttt{NumPy} optimizó el procesamiento de los cálculos estadísticos.

Esta base tecnológica permitió construir una herramienta de simulación de ruleta europea caracterizada por la aleatoriedad y la equiprobabilidad de sus resultados. Al contrastar las gráficas generadas con el marco teórico, se comprobó una alta correspondencia con los modelos matemáticos esperados. En particular, los estadísticos analizados tendieron a aproximarse a sus valores teóricos a medida que aumentó el número de tiradas.

Finalmente, la efectividad del programa quedó demostrada al evidenciar empíricamente la Ley de los Grandes Números y el Teorema Central del Límite. Se observó que, en la mayoría de las gráficas, cuando la cantidad de tiradas es suficientemente grande, las curvas tienden a sus respectivas asíntotas teóricas. Esto permite concluir que el simulador constituye un sistema válido y preciso para el análisis de fenómenos aleatorios simples, y que la ruleta europea resulta un caso adecuado para integrar programación, simulación y fundamentos estadísticos.

\newpage

\section{Referencias}

\begin{enumerate}[label={[\arabic*]}, leftmargin=*]

    \item Wikipedia. (2026). \textit{Ruleta}. 
    \url{https://es.wikipedia.org/wiki/Ruleta}

    \item Wikipedia. (2026). \textit{Probabilidad}. 
    \url{https://es.wikipedia.org/wiki/Probabilidad}

    \item Wikipedia. (2026). \textit{Estadística}. 
    \url{https://es.wikipedia.org/wiki/Estad%C3%ADstica}

    \item Wikipedia. (2026). \textit{Ley de los grandes números}. 
    \url{https://es.wikipedia.org/wiki/Ley_de_los_grandes_números}

    \item Khan Academy. (s. f.). \textit{Teorema Central del Límite}. 
    \url{https://es.khanacademy.org/math/ap-statistics/sampling-distributions-ap/what-is-the-central-limit-theorem/v/central-limit-theorem}

    \item Wikipedia. (2026). \textit{Método de Montecarlo}. 
    \url{https://es.wikipedia.org/wiki/Método_de_Montecarlo}

    \item Python Software Foundation. (s. f.). \textit{Módulo random — Generador de números pseudoaleatorios}. 
    \url{https://docs.python.org/es/3/library/random.html}

    \item Matplotlib Development Team. (s. f.). \textit{Matplotlib User Guide}. 
    \url{https://matplotlib.org/stable/users/index.html}

    \item NumPy Developers. (s. f.). \textit{NumPy Statistics Routines}. 
    \url{https://numpy.org/doc/stable/reference/routines.statistics.html}

\end{enumerate}

\end{document}
